\begin{explanationbox}
\textbf{一、先分清两个问题：斜率可以任意取，过定点后却只能取特定值。}

把直线想成一把可以转动、平移的直尺。
\begin{itemize}
\item \textbf{只规定方向，不规定经过哪里：}任取一个有限斜率 $k$，把直尺平移到刚好接触椭圆的位置，一定能找到切线；上下（沿法向的两个方向）各有一条。
\item \textbf{规定必须经过 $M$：}直尺不能自由平移，只能绕 $M$ 转动。若 $M$ 在椭圆外，恰好有两个临界方向能贴住椭圆；在椭圆内则不可能仅接触一次。
\item \textbf{直尺转到竖直：}此时斜率不存在，但切线仍然存在。斜率 $k$ 并不能描述所有方向。
\end{itemize}

\textbf{二、如何一眼记住 $m^2=a^2k^2+b^2$？}

把椭圆写成 $x=a\cos t,\ y=b\sin t$。在方向 $k$ 下观察
\[
y-kx=b\sin t-ak\cos t.
\]
其取值范围为
\[
-\sqrt{a^2k^2+b^2}\le y-kx\le\sqrt{a^2k^2+b^2}.
\]
直线 $y=kx+m$ 就是 $y-kx=m$：$m$ 在区间内部时通常穿过椭圆形成两个交点；只有位于两端点时才恰好相切。因此
\[
\boxed{\text{切线}\iff |m|=\sqrt{a^2k^2+b^2}.}
\]
\textbf{记忆口诀：}“先看斜率 $k$，再找截距极值；$m$ 平方等于 $a^2k^2+b^2$。”

\textbf{三、从切线判定式到过点斜率方程，只需一次代入。}

经过 $M$ 的斜线写为
\[
y=kx+m,\qquad m=y_0-kx_0.
\]
于是
\[
\underbrace{(y_0-kx_0)^2}_{\text{经过定点}}
=\underbrace{a^2k^2+b^2}_{\text{与椭圆相切}}.
\]
展开移项得到 (C009-1)。记住\textbf{“过点定截距，相切定平方”}，无需背诵整个二次方程。

\textbf{四、为什么外点两条、内点零条？}

把坐标按 $x\mapsto x/a,\ y\mapsto y/b$ 缩放，椭圆变成单位圆，点 $M$ 变成
$M'=(x_0/a,y_0/b)$。
缩放保持直线、相交及相切关系，也保持不同直线的条数。因此只要看
\[
OM'^2=S=\frac{x_0^2}{a^2}+\frac{y_0^2}{b^2}:
\quad S>1\text{ 外部},\ S=1\text{ 圆上},\ S<1\text{ 内部}.
\]
单位圆外一点能引两条切线，圆上只能引一条，圆内不能引切线。\textbf{注意：非等比例缩放不保持角度和垂直关系。}

\textbf{五、方程为什么有时从二次变成一次？}

当 $x_0=\pm a$，经过 $M$ 的竖直直线 $x=x_0$ 已经切于长轴顶点。由于 $y=kx+m$ 排除了竖直直线，剩下的方程只能找另一条斜切线；二次项消失并不意味着外点只有一条切线。

\textbf{六、数形联想与解题信号。}
\begin{itemize}
\item 题目出现“过定点作两条切线、求斜率和或积”时，优先用 (C009-1) 与韦达定理；先检查 $x_0^2\ne a^2$。
\item 题目出现“切线互相垂直、求切线交点轨迹”时，联想到蒙日圆 $x_0^2+y_0^2=a^2+b^2$；竖直与水平的特殊组合也要检查。
\item 题目出现“求切点或切点弦”时，切点弦为 $\frac{x_0x}{a^2}+\frac{y_0y}{b^2}=1$（外点引两条切线时，$x,y$ 表示两切点所在直线上的坐标）。
\end{itemize}
\end{explanationbox}
